\section{Related Work}
\label{sec:related}

\paragraph{Depth on mirrors, glass and displays.}
Mirrors are recovered as first-class scene elements with dedicated instance masks
and plane annotations, enabling depth to be refined on surfaces where standard
reconstruction fails \citep{mirror3d2021}, and transparency is modelled by
decomposing radiance into interface, transmitted and reflected components
\citep{glint2026}. Benchmarks quantify the difficulty directly, covering specular
and transparent surfaces with dense annotation \citep{booster2023}, large-scale
transparent objects \citep{clearpose2022}, and simulation-to-real depth
restoration for such materials \citep{dreds2022}. These establish that the
surfaces studied here are recognised as difficult and that recovering their
geometry is an active target. The present work takes a step back from recovering
the geometry and asks first whether the mechanism producing an appearance is
recoverable from the observations available.

\paragraph{Multi-layer and order-aware geometry.}
Representing more than one surface along a ray is well established. Ordinal
multi-layer benchmarks supply relations for both the first surface and the scene
behind it, together with real and synthetic training data
\citep{layereddepth2025}; self-determined grouping formulates transparency as
unordered events along a ray \citep{seegroup2026}; and steerable predictors
condition on a focus distance to select which layer is returned
\citep{depthfocus2026}. This line supplies the representational capacity that the
present formulation assumes, and one such benchmark is used here as an independent
evaluation set whose annotations share no machinery with the measures defined in
Section~\ref{sec:method}.

\paragraph{Apparent versus physical geometry.}
That predictors report depicted rather than physical structure is documented at
scale by a benchmark spanning illusion categories including printed pictures,
screens and mirrors \citep{illusion2025}. The distinction between apparent and
physical geometry is therefore taken as established. What is added here is a
criterion determining when that distinction is decidable from the observations an
observer can obtain, together with an explicit abstention when it is not, and a
benchmark constructed so that such undecidable cases are present and scorable.

\paragraph{Active perception and experimental design.}
Selecting viewpoints to maximise expected information is long established in
active and next-best-view perception, where the objective is typically a
belief-weighted sum over competing possibilities. The question posed here is
complementary: rather than which view is most informative, it asks whether any
view within a bounded action set is informative enough to license a commitment.
The criterion used for that purpose is not new. Designing experiments to
discriminate between rival models is a mature area of statistics, in which the
maximin criterion of \eqref{eq:maximin} is known as $T$-optimality
\citep{atkinson1975design,hunter1965experimental}; its role here is the transfer
of that criterion into active visual perception.

\paragraph{Causal identifiability, selective and conformal prediction.}
Identifiability under temporal interventions is established in causal
representation learning \citep{citris2022}; that interventional framing is
borrowed and applied to the narrower question of which optical mechanism produced
an image, with the action set playing the role of the intervention family.
Deferring on uncertain inputs follows Chow's rule \citep{chow1970}, and split
conformal prediction supplies distribution-free, finite-sample coverage guarantees
\citep{vovk2005,lei2018}. Those guarantees constrain how many predictions are
retained rather than which ones, so a calibrated procedure inherits the ranking of
the score it wraps. The abstention rule used here is derived instead from a
geometric condition on the available evidence, and conformal prediction is retained
as a baseline against which that condition is compared.
